% Headline numbers from paper/tables/scorecard.md and table1_bundle.tex.
Writers have no reliable way to show that a language model was trained on their work: membership-inference
attacks are close to random on large models when members and non-members come from the same distribution. We
study a creator-side alternative, \emph{canaries}: hidden, made-up details that a writer plants in their text and
later tests for. We ask how many copies of which canary a writer needs before training on their work becomes
provable at $p<0.01$, and with what access. We continue pretraining \pythia{160M} and \pythia{410M} on Wikipedia with four canary
types (random sequences, fluent trap sentences, invisible-character watermarks, and fictional facts) planted at
1, 4, 16 or 64 copies, with matched never-planted canaries as controls, and pass the corpus through realistic
cleaning steps (substring deduplication, Unicode normalization, paraphrasing). A pre-registered one-sided
Mann--Whitney test on member vs.\ non-member scores gives exact false-positive control, which we verify on the
untouched base models (0 of 32 tests significant at each size). With log-probability access, the larger model
detects every canary type from a single copy, but only when repeated canaries of the same kind are also in the
corpus: once deduplication removes those repeats, the same single copies are barely detectable. Unicode
normalization strips watermarks, and substring deduplication removes every verbatim canary type. Only fictional
facts, phrased differently in each copy, survive all three cleaning steps (16 copies for a 25-fact bundle on
\pythia{410M}). With text-only access, verbatim canaries become detectable from 4 copies on the larger model, but
fictional facts are never reliably detected.
